\textbf{Phần 1. Trắc nghiệm lựa chọn}
\setcounterpageref{ex}{0}

\begin{ex} %[3D3V68]
Bạn Mai có 2 tờ tiền loại $20\,000$ đồng và 1 tờ loại $10\,000$ đồng. Tổng số tiền Mai có là
\choice
{$40\,000$ đồng}
{\True $50\,000$ đồng}
{$30\,000$ đồng}
{$60\,000$ đồng}
\loigiai{
$20\,000 \times 2 + 10\,000 = 50\,000$ đồng.
\textbf{Đáp án của câu hỏi là B}
}
\end{ex}

\begin{ex} %[3D3V68]
Một chiếc bút mực giá $35\,000$ đồng. Bạn Nam đưa cô bán hàng tờ tiền $50\,000$ đồng. Cô bán hàng phải trả lại Nam bao nhiêu tiền thừa?
\choice
{$25\,000$ đồng}
{$10\,000$ đồng}
{$20\,000$ đồng}
{\True $15\,000$ đồng}
\loigiai{
$50\,000 - 35\,000 = 15\,000$ đồng.
\textbf{Đáp án của câu hỏi là D}
}
\end{ex}

\begin{ex} %[3D3V68]
Đổi một tờ tiền $100\,000$ đồng lấy các tờ tiền mệnh giá $20\,000$ đồng. Số tờ tiền $20\,000$ đồng nhận được là
\choice
{\True $5$ tờ}
{$4$ tờ}
{$6$ tờ}
{$10$ tờ}
\loigiai{
$100\,000 : 20\,000 = 5$ tờ.
\textbf{Đáp án của câu hỏi là A}
}
\end{ex}

\begin{ex} %[3D3V68]
Một quyển truyện tranh giá $45\,000$ đồng, một hộp bút màu giá $38\,000$ đồng. Mua cả hai món đồ hết tất cả bao nhiêu tiền?
\choice
{$73\,000$ đồng}
{$85\,000$ đồng}
{\True $83\,000$ đồng}
{$80\,000$ đồng}
\loigiai{
$45\,000 + 38\,000 = 83\,000$ đồng.
\textbf{Đáp án của câu hỏi là C}
}
\end{ex}

\begin{ex} %[3D3V68]
Để có đúng $50\,000$ đồng bằng các tờ tiền loại $5000$ đồng thì cần bao nhiêu tờ tiền?
\choice
{$5$ tờ}
{\True $10$ tờ}
{$20$ tờ}
{$15$ tờ}
\loigiai{
$50\,000 : 5000 = 10$ tờ.
\textbf{Đáp án của câu hỏi là B}
}
\end{ex}

\begin{ex} %[3D3V68]
An mua 4 quyển vở, mỗi quyển giá $12\,000$ đồng và 1 chiếc thước kẻ giá $5000$ đồng. An đưa cô bán hàng tờ tiền $100\,000$ đồng. Cô bán hàng trả lại An số tiền thừa là
\choice
{$53\,000$ đồng}
{$45\,000$ đồng}
{$48\,000$ đồng}
{\True $47\,000$ đồng}
\loigiai{
Tổng tiền mua: $12\,000 \times 4 + 5000 = 48\,000 + 5000 = 53\,000$ đồng. Tiền thừa: $100\,000 - 53\,000 = 47\,000$ đồng.
\textbf{Đáp án của câu hỏi là D}
}
\end{ex}

\begin{ex} %[3D3V68]
Mệnh giá tiền giấy polymer nào sau đây KHÔNG tồn tại trong hệ thống tiền tệ hiện hành của Việt Nam?
\choice
{\True $30\,000$ đồng}
{$20\,000$ đồng}
{$50\,000$ đồng}
{$100\,000$ đồng}
\loigiai{
Không có tờ tiền mệnh giá $30\,000$ đồng.
\textbf{Đáp án của câu hỏi là A}
}
\end{ex}

\begin{ex} %[3D3V68]
Bình có 3 tờ $10\,000$ đồng và 4 tờ $5000$ đồng. Bình muốn mua một quả bóng giá $60\,000$ đồng. Hỏi Bình còn thiếu bao nhiêu tiền?
\choice
{$15\,000$ đồng}
{$20\,000$ đồng}
{\True $10\,000$ đồng}
{$5000$ đồng}
\loigiai{
Bình có: $10\,000 \times 3 + 5000 \times 4 = 30\,000 + 20\,000 = 50\,000$ đồng. Còn thiếu: $60\,000 - 50\,000 = 10\,000$ đồng.
\textbf{Đáp án của câu hỏi là C}
}
\end{ex}

\begin{ex} %[3D3V68]
Một chiếc cặp sách giá $85\,000$ đồng. Nhân dịp khai giảng được giảm giá $15\,000$ đồng. Giá tiền của chiếc cặp sách sau khi giảm là
\choice
{$75\,000$ đồng}
{\True $70\,000$ đồng}
{$80\,000$ đồng}
{$65\,000$ đồng}
\loigiai{
$85\,000 - 15\,000 = 70\,000$ đồng.
\textbf{Đáp án của câu hỏi là B}
}
\end{ex}

\begin{ex} %[3D3V68]
Có bao nhiêu cách đổi tờ tiền $20\,000$ đồng ra các tờ tiền mệnh giá $5000$ đồng và $10\,000$ đồng (dùng cả hai loại tiền)?
\choice
{$2$ cách}
{$3$ cách}
{Không đổi được}
{\True $1$ cách (gồm 1 tờ $10\,000$ đồng và 2 tờ $5000$ đồng)}
\loigiai{
Để có cả 2 loại tiền: $20\,000 = 10\,000 \times 1 + 5000 \times 2$. Chỉ có duy nhất 1 cách.
\textbf{Đáp án của câu hỏi là D}
}
\end{ex}

\begin{ex} %[3D3V68]
Mẹ đi siêu thị mua một khay thịt gà hết $65\,000$ đồng và một bó rau hết $15\,000$ đồng. Mẹ đưa 2 tờ tiền loại $50\,000$ đồng. Số tiền thừa nhận lại là
\choice
{\True $20\,000$ đồng}
{$30\,000$ đồng}
{$10\,000$ đồng}
{$25\,000$ đồng}
\loigiai{
Mẹ đưa: $50\,000 \times 2 = 100\,000$ đồng. Mua hết: $65\,000 + 15\,000 = 80\,000$ đồng. Tiền thừa: $100\,000 - 80\,000 = 20\,000$ đồng.
\textbf{Đáp án của câu hỏi là A}
}
\end{ex}

\begin{ex} %[3D3V68]
Mua 5 chiếc bánh bao hết $50\,000$ đồng. Hỏi mua 8 chiếc bánh bao như thế hết bao nhiêu tiền?
\choice
{$75\,000$ đồng}
{$85\,000$ đồng}
{\True $80\,000$ đồng}
{$70\,000$ đồng}
\loigiai{
Một chiếc bánh bao: $50\,000 : 5 = 10\,000$ đồng. Tám chiếc bánh bao: $10\,000 \times 8 = 80\,000$ đồng.
\textbf{Đáp án của câu hỏi là C}
}
\end{ex}

\textbf{Phần 2. Trắc nghiệm đúng sai}
\setcounterpageref{ex}{0}

\begin{ex} %[3D3V68]
Xét tính đúng sai của các nhận định về tiền tệ Việt Nam:
\choiceTF
{\True Đơn vị tiền tệ của nước ta là đồng.}
{Một tờ 50 000 đồng đổi được đúng 10 tờ 10 000 đồng.}
{\True Hai tờ 50 000 đồng có giá trị bằng một tờ 100 000 đồng.}
{Tờ tiền 100 000 đồng có giá trị nhỏ hơn tờ tiền 50 000 đồng.}
\loigiai{
\begin{itemize}
    \item Ý a) Đúng: Quy chuẩn quốc gia.
    \item Ý b) Sai: Sai vì $50\,000 : 10\,000 = 5$ tờ.
    \item Ý c) Đúng: $50\,000 \times 2 = 100\,000$ đồng.
    \item Ý d) Sai: Sai vì $100\,000 > 50\,000$.
\end{itemize}
\begin{itemize}
    \item a) \textbf{Đúng}.
    \item b) \textbf{Sai}.
    \item c) \textbf{Đúng}.
    \item d) \textbf{Sai}.
\end{itemize}
}
\end{ex}

\begin{ex} %[3D3V68]
Bạn Cúc có một con heo đất tiết kiệm gồm: 5 tờ 10 000 đồng, 4 tờ 5000 đồng và 10 tờ 2000 đồng:
\choiceTF
{\True Tổng số tiền loại 10 000 đồng là 50 000 đồng.}
{\True Tổng số tiền loại 5000 đồng là 20 000 đồng.}
{\True Tổng số tiền loại 2000 đồng là 20 000 đồng.}
{Cúc có tất cả đúng 100 000 đồng trong heo đất.}
\loigiai{
\begin{itemize}
    \item Ý a) Đúng: $10\,000 \times 5 = 50\,000$ đồng.
    \item Ý b) Đúng: $5000 \times 4 = 20\,000$ đồng.
    \item Ý c) Đúng: $2000 \times 10 = 20\,000$ đồng.
    \item Ý d) Sai: Sai vì tổng là $50 + 20 + 20 = 90\,000$ đồng.
\end{itemize}
\begin{itemize}
    \item a) \textbf{Đúng}.
    \item b) \textbf{Đúng}.
    \item c) \textbf{Đúng}.
    \item d) \textbf{Sai}.
\end{itemize}
}
\end{ex}

\textbf{Phần 3. Trả lời ngắn}
\setcounterpageref{ex}{0}

\begin{ex} %[3D3V68]
Một tờ 100 000 đồng đổi được bao nhiêu tờ 50 000 đồng?
\shortans{2}
\loigiai{
\textbf{Đáp số:} $2$\\
$100\,000 : 50\,000 = 2$ tờ.
}
\end{ex}

\begin{ex} %[3D3V68]
Lan có 3 tờ 20 000 đồng và 2 tờ 10 000 đồng. Hỏi Lan có tất cả bao nhiêu nghìn đồng?
\shortans{80}
\loigiai{
\textbf{Đáp số:} $80$\\
$20 \times 3 + 10 \times 2 = 60 + 20 = 80$ nghìn đồng.
}
\end{ex}

\begin{ex} %[3D3V68]
Mua 3 cốc trà sữa giá 25 000 đồng mỗi cốc. Hỏi phải trả tất cả bao nhiêu nghìn đồng?
\shortans{75}
\loigiai{
\textbf{Đáp số:} $75$\\
$25 \times 3 = 75$ nghìn đồng.
}
\end{ex}

\begin{ex} %[3D3V68]
Bác Nam mua đồ hết 72 000 đồng và đưa tờ 100 000 đồng. Người bán hàng trả lại bác bao nhiêu nghìn đồng?
\shortans{28}
\loigiai{
\textbf{Đáp số:} $28$\\
$100 - 72 = 28$ nghìn đồng.
}
\end{ex}

\begin{ex} %[3D3V68]
Cần bao nhiêu tờ 20 000 đồng để có được đúng 100 000 đồng?
\shortans{5}
\loigiai{
\textbf{Đáp số:} $5$\\
$100\,000 : 20\,000 = 5$ tờ.
}
\end{ex}

\begin{ex} %[3D3V68]
Một cuốn sổ giá 18 000 đồng, một chiếc bút giá 7000 đồng. Mua cả hai thứ hết bao nhiêu nghìn đồng?
\shortans{25}
\loigiai{
\textbf{Đáp số:} $25$\\
$18 + 7 = 25$ nghìn đồng.
}
\end{ex}
